\documentclass{leveldoc}
\DocUses{foundations,first_bend,slingshot,injection}
\DocFingerprint{b6910f86734d5696}

\begin{document}

\section{What the level asks}\label{sec:asks}

The game's text:
\leveldescription

One particle ($q = \val{q}$, $m = \val{m}$) is launched from $(\val{xLaunch}, \val{yLaunch})$
along $+x$, with kinetic energy $T_0 = \val{Tzero}$. Nine equal charges stand in a wall at
$x = \val{xWall}$, from $y = \val{yWallLo}$ to $y = \val{yWallHi}$, and block that line.
The detector is the rectangle $\val{xDlo} \le x \le \val{xDhi}$,
$\val{yDlo} \le y \le \val{yDhi}$. An entry counts only when the velocity there lies within
$\val{halfDeg}^\circ$ of $+x$. The player places at most $\val{maxCharges}$ charges, of
either sign, from the magnitudes in the level file, and no magnets. A flight ends when the
particle touches a charge, leaves the world, enters the detector or runs out of time
(\xsecref{foundations}{sec:board}); $t_{\max} = \val{tmax}$.

The reference solution adds $Q = \val{Qup}$ at $(\val{xUp}, \val{yUp})$ and
$Q = \val{Qat}$ at $(\val{xAt}, \val{yAt})$. The first stands just ahead of and above the
launch and pushes the beam down. The second sits below the wall and attracts, so the beam
passes under the wall and climbs back. It meets the detector on the left face, at
$\val{angE}^\circ$ to $+x$, inside the cone. This document integrates the motion of the
eleven charges and applies the cone of \levelref{injection}.

\section{Setup}\label{sec:setup}

The level file, field by field (the appendix reproduces it):
\setuptable

In symbols: $c = \val{c}$, the particle has radius $a = \val{aRad}$, and every charge has
radius $R = \val{R}$. The launch has
\[
  \gamma_0 = 1 + \frac{T_0}{mc^2} = \val{gammaZero}, \qquad
  \abs{\vect{p}_0} = \sqrt{\frac{T_0^2}{c^2} + 2mT_0} = \val{pZero}, \qquad
  v_0 = \frac{\abs{\vect{p}_0}}{\gamma_0 m} = \val{vZero}.
\]
The direction in the file is $+x$. The world is $-1 \le x \le \val{xWorld}$,
$-1 \le y \le \val{yWorld}$. The wall is $\val{nWall}$ charges of value $\val{Qwall}$, so
each coupling is $\kappa = qQ = \val{kWall}$. The two added couplings are
$\kappa_\mathrm{u} = \val{kUp}$ and $\kappa_\mathrm{a} = \val{kAt}$. The cone's axis is
the angle $\val{axisDeg}^\circ$ and its half-angle is $\val{halfDeg}^\circ$, so the axis
is the unit vector $\uvec{e} = (1, 0)$ and the half-angle in radians is
$\delta = \val{halfRad}$. The detector sets no condition on the kinetic energy.

\begin{figure}[htbp]
\centering
\begin{tikzpicture}
\begin{axis}[axis equal image, width=0.88\linewidth, xmin=-1.6, xmax=32.6, ymin=-2.2, ymax=21.6,
  xtick={0,5,...,30}, ytick={0,5,...,20}, tick label style={font=\small},
  xlabel={$x$}, ylabel={$y$}, label style={font=\small}]
  \draw[black!45, dashed] (-1,-1) rectangle (31,21);
  \draw[black!50] (0,0) rectangle (30,20);
  \fill[green!50!black, fill opacity=0.16] (28,8) rectangle (30,12);
  \draw[green!40!black] (28,8) rectangle (30,12);
  \fill[teal, fill opacity=0.30] (\rawval{coneOx},\rawval{coneOy}) --
    (\rawval{coneXp},\rawval{coneYp}) -- (\rawval{coneXm},\rawval{coneYm}) -- cycle;
  \addplot[only marks, mark=*, mark size=1.15pt, red!75!black] table {wall.dat};
  \fill[blue!70!black, fill opacity=0.35] (15,1) circle[radius=0.3];
  \draw[blue!70!black] (15,1) circle[radius=0.3];
  \fill[red!70!black] (2,13) circle[radius=0.22];
  \addplot[thin, black!55] table {pathup.dat};
  \addplot[thick, orange!80!black] table {pathat.dat};
  \addplot[very thick, blue!70!black] table {path.dat};
  \draw[-{Stealth[length=6pt]}, thick] (\rawval{xLaunch},\rawval{yLaunch}) --
    (\rawval{arrowX},\rawval{arrowY});
  \fill (\rawval{xLaunch},\rawval{yLaunch}) circle[radius=0.11];
  \node[font=\scriptsize, red!60!black, anchor=south] at (2,13.35) {$Q_\mathrm{u}$};
  \node[font=\scriptsize, blue!50!black, anchor=west] at (15.45,1) {$Q_\mathrm{a}$};
  \node[font=\scriptsize, red!60!black, anchor=south] at (15,14.35) {wall};
  \node[font=\scriptsize, green!30!black, anchor=south] at (29,12.15) {detector};
  \node[font=\scriptsize, teal!80!black, anchor=west] at (30.55,10) {$\pm 8^\circ$};
\end{axis}
\end{tikzpicture}
\caption{The reference flight (blue) passes under the wall and the attracting charge, then
enters the left face of the detector. The teal wedge is the cone of accepted directions,
$\pm\val{halfDeg}^\circ$ about $+x$, drawn from the middle of that face. The upper charge
alone (grey) and the attracting charge alone (orange) both leave through the floor of the
world. The circle on $Q_\mathrm{a}$ has radius $R$.}
\label{fig:board}
\end{figure}

\section{Principles and model}\label{sec:model}

\begin{principles}
  \item The field of a charged sphere and the force on a non-polarizable particle,
    \xeqref{foundations}{eq:coulomb} and \xeqref{foundations}{eq:sphereforce}.
  \item The relativistic equation of motion \xeqref{foundations}{eq:eom}, the launch
    \xeqref{foundations}{eq:launch} and the energy \xeqref{foundations}{eq:energy}. The
    motion stays in the plane (\xsecref{foundations}{sec:plane}).
  \item The board's rules (\xsecref{foundations}{sec:board}). The first entry into the
    detector decides, and the velocity there has to lie in the cone
    \xeqref{injection}{eq:cone}.
  \item The classical Runge--Kutta method \xeqref{foundations}{eq:rk4} and the fixed-step
    flight of \xsecref{slingshot}{sec:method}. With two or more centres the orbit equation
    does not separate (\xeqref{slingshot}{eq:noseparation}).
  \item The radiated power \xeqref{foundations}{eq:larmor}, with the acceleration
    \xeqref{first_bend}{eq:acceleration}, summed by Simpson's rule
    \xeqref{foundations}{eq:simpson}.
\end{principles}

The level has no magnets, so $\vect{B} = \vect{0}$. With the charges at
$\vect{r}_i$ and couplings $\kappa_i = qQ_i$, in the file's order (the wall, then the
charges the player added), the force is
\begin{equation}\label{eq:force}
  \vect{F}(\vect{x}) = \sum_{i=1}^{N} \kappa_i \frac{\vect{x} - \vect{r}_i}{\abs{\vect{x} - \vect{r}_i}^3},
  \qquad N = \val{nCharges}
\end{equation}
on the reference placement. The state $(x, y, p_x, p_y)$ obeys \xeqref{foundations}{eq:eom}
with this force until the first event. The same superposition gives the potential energy
$V$ in
\begin{equation}\label{eq:energy}
  W = \gamma mc^2 + V(\vect{x}) = \gamma_0 mc^2 + V(\vect{x}_0).
\end{equation}

\begin{assumption}{The charges are held in place}\label{as:fixed}
None of the charges recoils. That is a rule of the game
(\xsecref{foundations}{sec:radiation}).
\end{assumption}

\begin{assumption}{No radiation}\label{as:radiation}
The radiated energy is not taken from the particle. Along the reference flight it is
$\val{Erad}$, the fraction $\val{radFrac}$ of $T_0$ (§\ref{sec:hand}). The two flights of
§\ref{sec:others} radiate $\val{UpRad}$ and $\val{AtRad}$. The game makes the same
approximation and reports the same integral (\physref{8}).
\end{assumption}

\begin{assumption}{A point particle}\label{as:point}
$a = 0$, and the particle does not polarize, so \xeqref{foundations}{eq:sphereforce} is
exact. On the reference flight the closest approach to the attracting charge is
$\val{rClose}$, which is $\val{rOverR}$ times $R$, a clearance of $\val{clearAt}$. The
nearest charge of the wall stays $\val{rWall}$ away, a clearance of $\val{clearWall}$.
\end{assumption}

\begin{assumption}{The step resolves the close pass}\label{as:step}
The reported step is $h = \val{H}$. Halving it estimates the error of the reference end
state as $\val{errRef}$ (§\ref{sec:hand}). That is far below the half-angle
$\delta = \val{halfRad}$ and far below the height of the detector. The tolerance used in
§\ref{sec:comparison} is ten times this estimate.
\end{assumption}

\section{Derivation}\label{sec:derivation}

\subsection{What the energy allows}

At the launch the upper charge is a distance $\val{rUp}$ away and the attracting charge a
distance $\val{rAt}$, so $V(\vect{x}_0) = \val{VZero}$ and
\[
  W - mc^2 = T_0 + V(\vect{x}_0) = \val{Wbind}.
\]
This is positive. Far from every charge $V \to 0$, which is below the bound, so
\eqref{eq:energy} does not confine the particle: leaving the world is allowed, and the
integration has to decide whether it does. The force at the launch is
$(\val{FxZero}, \val{FyZero})$. Both player charges point it downward. The upper one
contributes $\val{FyUp}$ and the attracting one $\val{FyAtLaunch}$.

\subsection{Why the flight is integrated}\label{sec:noseparation}

With one charge, angular momentum closes the motion into the orbit of
\levelref{first_bend}. Two fixed centres already fail to separate in the relativistic
problem: the term \xeqref{slingshot}{eq:noseparation} does not split. Eleven centres give
that obstruction no symmetry that would remove it. The flight is the initial-value problem
for \eqref{eq:force}, and §\ref{sec:method} computes it to a stated accuracy.

\subsection{The direction cone}\label{sec:cone}

The test is the one derived in \levelref{injection}. The level file stores the cone as an
axis $\alpha = \val{axisDeg}^\circ$ and a half-angle $\delta = \val{halfDeg}^\circ$. The
axis is $\uvec{e} = (1, 0)$. If $\vect{v}$ is the velocity at the first entry into the box,
\begin{equation}\label{eq:cone}
  \beta = \operatorname{atan2}\bigl(\abs{\vect{v}\times\uvec{e}},\, \vect{v}\cdot\uvec{e}\bigr),
\end{equation}
which is \xeqref{injection}{eq:cone}, and the entry counts when $\beta \le \delta$ with
$\delta$ in radians. The outcome is then Arrived. When the particle has entered and
$\beta > \delta$, the box absorbs it and the entry does not count: the outcome is Rejected
(\physref{6.1}). A flight that never enters the box is not tested. Since
$\vect{v} = \vect{p}/(\gamma m)$ with $\gamma m > 0$, the same $\beta$ comes from
$\vect{p}$. Here that is $\beta = \operatorname{atan2}(\abs{p_y}, p_x)$.

\subsection{The method}\label{sec:method}

The state is advanced by the classical Runge--Kutta method \xeqref{foundations}{eq:rk4} at
the fixed step $h$, as in \xsecref{slingshot}{sec:method}. After each step the position on
that step is the cubic Hermite interpolant \xeqref{slingshot}{eq:hermite}, sampled at
$\val{nSample}$ points. The speed is less than $c$, and the velocity terms of the
interpolant add less than $ch$, so the sample stays inside the box of the step's two ends
enlarged by $ch$. A sphere, the world's edge or the detector is tested only when that box
meets it. The samples are read in the board's order
(\xsecref{foundations}{sec:board}): a sphere of radius $R + a$ about a charge, then the
world's edge, then the detector. The earliest sample wins, and a tie keeps that order. The
first instant of the winning event is then the bisection of the interpolant,
$\val{nBisect}$ iterations. The position reported is the interpolant at that instant. The
momentum is one further step of the method, of that length.

When the event is a face of a box and the free coordinate lies in that face, the face
coordinate is set to the face value. The residual of the interpolant, before that
assignment, is the check that the step really did cross the face. The error of the
reported step is a fifteenth of its difference from the step $2h$
(\xsecref{foundations}{sec:numerics}). The energy \eqref{eq:energy} is a second check, one
that does not use the game.

The closest approach quoted below is the minimum of the same interpolant. On a step whose
start lies more than $2ch$ beyond the minimum already found, that charge is skipped: the
distance to it cannot fall by more than that.

\section{The reference solution by hand}\label{sec:hand}

\begin{procedure}
  \item\label{st:launch} \emph{The launch} (\xeqref{foundations}{eq:launch}):
    $\gamma_0 = \val{gammaZero}$ and $\abs{\vect{p}_0} = \val{pZero}$, aimed along $+x$.
    The force is $(\val{FxZero}, \val{FyZero})$, so the beam starts downward.
  \item\label{st:energy} \emph{The energy} \eqref{eq:energy}: $W = \val{W}$ and
    $W - mc^2 = \val{Wbind} > 0$. The world's edge is energetically open.
  \item\label{st:flight} \emph{The flight}: the method of §\ref{sec:method} with
    $h = \val{H}$. The states every five time units, and the drift of \eqref{eq:energy}:
    \begin{center}\small\val{states}\end{center}
    At $t = \val{tFive}$ the particle is at $(\val{xFive}, \val{yFive})$, with
    $p_y = \val{pyFive}$. At $t = \val{tFifteen}$ it is at
    $(\val{xFifteen}, \val{yFifteen})$ and $p_y = \val{pyFifteen}$ is still negative.
    The lowest point is $(\val{xDip}, \val{yDip})$ at $t = \val{tDip}$. The path crosses
    the wall's line $x = \val{xWall}$ at $y = \val{yPass}$, below the wall, and the
    nearest wall charge (the one at $y = \val{yWallHit}$) is $\val{rWall}$ away.
  \item\label{st:dip} \emph{Under the attracting charge}: the closest approach is
    $\val{rClose}$, at $(\val{xClose}, \val{yClose})$ and $t = \val{tClose}$, a clearance
    $\val{clearAt}$ outside the sphere. The force there is
    $(\val{FxClose}, \val{FyClose})$, upward. At $x = \val{xRise}$ the height is
    $\val{yRise}$. By $t = \val{tTwentyFive}$ the particle has climbed to
    $(\val{xTwentyFive}, \val{yTwentyFive})$ with $p_y = \val{pyTwentyFive} > 0$.
    At $t = \val{tThirty}$ it is at $(\val{xThirty}, \val{yThirty})$, still climbing,
    with $p_y = \val{pyThirty}$.
  \item\label{st:entry} \emph{The entry}: the step from $t_n = \val{tBefore}$, at
    $(\val{xBefore}, \val{yBefore})$, ends at $(\val{xAfter}, \val{yAfter})$. The
    interpolant crosses the \val{faceE} face at the fraction $u = \val{uE}$. The residual
    before the face coordinate is assigned is $\val{faceRes}$, so
    $x_\mathrm{e} = \val{xE}$, $y_\mathrm{e} = \val{yE}$ and
    $t_\mathrm{e} = \val{tE}$, after $\val{steps}$ steps. One step of length
    $t_\mathrm{e} - t_n$ gives $\vect{p}_\mathrm{e} = (\val{pxE}, \val{pyE})$.
  \item\label{st:cone} \emph{The cone} \eqref{eq:cone}:
    $\beta = \val{betaE}$ radians, which is $\val{angE}^\circ$. The margin under
    $\val{halfDeg}^\circ$ is $\val{marginDeg}^\circ$, so the entry counts.
  \item\label{st:check} \emph{The checks}: the energy gives
    $\gamma_\mathrm{e} = (W - V(\vect{x}_\mathrm{e}))/(mc^2) = \val{gammaE}$ and
    $\abs{\vect{p}_\mathrm{e}} = mc\sqrt{\gamma_\mathrm{e}^2 - 1} = \val{pEnergy}$, which
    the integration matches to $\val{pDiff}$. Along the flight the energy stays within
    $\val{wDev}$ of $W$. Halving the step estimates the error of the end state as
    $\val{errRef}$. Coarser steps cross the same face:
    \begin{center}\small\val{convergence}\end{center}
  \item\label{st:rad} \emph{The radiation}: the power \xeqref{foundations}{eq:larmor} with
    the acceleration \xeqref{first_bend}{eq:acceleration}, summed by Simpson's rule on the
    steps of $h/2$. The integral is $\val{Erad}$, the fraction $\val{radFrac}$ of $T_0$.
    The relative difference from the same sum at step $h$, divided by fifteen, is
    $\val{radErr}$.
\end{procedure}

\begin{result}[The reference flight]
The particle enters the detector through its left face at $(\val{xE}, \val{yE})$ at
$t_\mathrm{e} = \val{tE}$, with momentum $(\val{pxE}, \val{pyE})$. The direction is
$\val{angE}^\circ$ from $+x$, inside the cone by $\val{marginDeg}^\circ$. The level is
solved.
\end{result}

By hand, Step~\ref{st:flight} is $\val{steps}$ steps of four evaluations of \eqref{eq:force}.
The close pass at a distance $\val{rClose}$ is what sets the step: the error estimate
$\val{errRef}$ is the reference flight's, and the two smoother flights of §\ref{sec:others}
come out smaller.

\section{Comparison with the game}\label{sec:comparison}

\comparisontable

The three outcomes agree. The tolerance on a state is $\val{tolFlight}$: ten times
$\val{errRef}$, the largest of the three step-halving estimates, and at least $10^{-9}$.
It was fixed from those estimates before the comparison. The largest absolute difference
is the reference ordinate, $\val{cmpRefYDiff}$. The estimate for that flight is
$\val{errRef}$, and the game's verified positions are good to about $10^{-10}$
(\xsecref{foundations}{sec:verification}), so a difference of this size is what the two
errors allow. The reference abscissa differs by $\val{cmpRefXDiff}$. The two floor
ordinates differ by $\val{cmpUpYDiff}$ and $\val{cmpAtYDiff}$.

The single-charge estimates are $\val{UpErr}$ and $\val{AtErr}$. Their absolute
differences stay smaller than the reference ordinate's. The radiated energies are compared
relatively, with tolerance $10^{-4}$, set the same way and wider than the quadrature
uncertainties $\val{radErr}$, $\val{UpRadErr}$ and $\val{AtRadErr}$. The largest relative
difference is the attracting charge's, $\val{cmpAtRadDiff}$.

\section{Other solutions}\label{sec:others}

Either added charge, without the other, leaves the world through the floor. The method of
§\ref{sec:method} integrates both (\code{placements.json}), each at $h = \val{H}$ and at
$2h$.
\begin{enumerate}[label=\arabic*.]
  \item $Q = \val{Qup}$ at $(\val{xUp}, \val{yUp})$ only. The particle leaves through the
    \val{UpFace} of the world at $(\val{UpX}, \val{UpY})$, at $t = \val{UpT}$.

    The momentum is $(\val{UpPx}, \val{UpPy})$. Its horizontal part is negative, and the
    exit is still to the left of the wall: the charge has turned the beam back, and the
    beam falls out of the board. The step-halving estimate is $\val{UpErr}$. The radiated
    energy is $\val{UpRad}$.
  \item $Q = \val{Qat}$ at $(\val{xAt}, \val{yAt})$ only. The particle leaves through the
    \val{AtFace} at $(\val{AtX}, \val{AtY})$, at $t = \val{AtT}$.

    The momentum is $(\val{AtPx}, \val{AtPy})$. Its horizontal part is still positive, and
    the exit is still to the left of the wall: the charge pulls the beam through the floor
    before the beam can pass under the wall. The estimate is $\val{AtErr}$. The radiated
    energy is $\val{AtRad}$.
\end{enumerate}

The reference needs both charges. The upper one supplies the downward push at launch
(Step~\ref{st:launch}); the attracting one supplies the upward force at the bottom of the
dip (Step~\ref{st:dip}), which is what brings the beam back up to the detector and into
the cone. Figure~\ref{fig:board} shows the three paths.

\section{What the game computes}\label{sec:engine}

The game integrates \xeqref{foundations}{eq:eom} with its eighth-order Runge--Kutta method
at the tolerances $10^{-10}$ and $10^{-12}$ (\physref{5}). It finds the crossing on its
interpolant (\physref{6}) and applies \eqref{eq:cone} to the velocity at that entry
(\physref{6.1}). It verifies the outcome by comparing its two runs (\physref{7}). The
model is this document's, including \ref{as:radiation} (\physref{8}). The agreement in
§\ref{sec:comparison} is between the fixed-step flight of §\ref{sec:method}, the cone
\eqref{eq:cone} and that integration.

\section{Further reading}\label{sec:reading}

\begin{itemize}
  \item \levelref{slingshot}: the fixed-step method, the Hermite crossing, and the reason
    two fixed centres have no closed relativistic orbit.
  \item \levelref{injection}: the cone \eqref{eq:cone}. There a single charge bends one
    ray into the same $\pm 8^\circ$ of $+x$.
  \item \levelref{first_bend}: the acceleration that enters the radiated power.
  \item \xsecref{foundations}{sec:board} and \xsecref{foundations}{sec:numerics}: the
    order of events, and the step-halving estimate.
\end{itemize}

\levelfileappendix

\end{document}
